\documentclass[final]{beamer}
\usepackage[size=a0,orientation=portrait,scale=1.13]{beamerposter}
\usetheme{gemini}
\usecolortheme{csl2s}
\usepackage[english]{babel}
\usepackage{newtxtext,newtxmath}
\usepackage{amsmath,amssymb}
\usepackage{graphicx}
\usepackage{booktabs,tabularx,array}
\usepackage{tikz}
\usetikzlibrary{arrows.meta,positioning,calc}
\usepackage{pgfplots}
\pgfplotsset{compat=1.18}
\setbeamertemplate{navigation symbols}{}
\setbeamertemplate{caption}[numbered]
\setbeamerfont{headline title}{family=\Raleway,size=\fontsize{72}{80}\selectfont,series=\bfseries}
\setbeamerfont{headline author}{family=\Raleway,size=\fontsize{25}{30}\selectfont}
\setbeamerfont{headline institute}{family=\Raleway,size=\fontsize{21}{26}\selectfont}
\setbeamerfont{block title}{family=\Raleway,size=\fontsize{40}{46}\selectfont,series=\bfseries}
\setbeamerfont{block body}{family=\rmfamily,size=\fontsize{28}{34}\selectfont}
\setbeamerfont{caption}{family=\rmfamily,size=\fontsize{21}{24}\selectfont}
\setbeamerfont{itemize/enumerate body}{family=\rmfamily,size=\fontsize{26}{31}\selectfont}
\setbeamerfont{itemize/enumerate subbody}{family=\rmfamily,size=\fontsize{24}{29}\selectfont}
\newlength{\sepwidth}
\newlength{\colwidth}
\newlength{\widecolwidth}
\setlength{\sepwidth}{0.03\paperwidth}
\setlength{\colwidth}{0.29\paperwidth}
\setlength{\widecolwidth}{0.61\paperwidth}
\title{Can ten qubits finish your sentence?}
\author{Ethan Bandasack\samelineand Malo Leroy\samelineand Thierry Zhang}
\institute{Not better than a small classical network: three quantum language models tested against fair controls.\\[0.25em]
CentraleSupélec / L2S \quad Supervisor: Zeno Toffano}
\logoleft{\colorbox{white}{\includegraphics[width=14cm]{logos/CS.png}}}
\logoright{\colorbox{white}{\includegraphics[width=13cm]{logos/L2S_logo.png}}}
\pgfplotsset{
  every axis/.append style={
    axis line style={line width=3pt,color=navy},
    tick style={line width=3pt,color=navy},
    grid=major,
    grid style={line width=1pt,color=pale},
    label style={font=\fontsize{24}{28}\selectfont},
    tick label style={font=\fontsize{24}{28}\selectfont},
    legend style={font=\fontsize{24}{28}\selectfont,draw=none,fill=none},
    mark size=5pt,
    clip=false
  }
}
\begin{document}
\begin{frame}[t]
\begin{columns}[t,totalwidth=\textwidth]
  \begin{column}{\colwidth}
    \begin{block}{Meaning as operators}
      Quantum semantics relates meanings.
    \end{block}
  \end{column}
  \begin{column}{\sepwidth}\end{column}
  \begin{column}{\colwidth}
    \begin{block}{Thin evidence}
      Metrics may obscure word recovery.
    \end{block}
  \end{column}
  \begin{column}{\sepwidth}\end{column}
  \begin{column}{\colwidth}
    \begin{block}{Our test}
      We test matched classical controls.
    \end{block}
  \end{column}
\end{columns}

\vspace{1cm plus 1fill}

\begin{columns}[t,totalwidth=\textwidth]
  \begin{column}{\colwidth}
    \begin{block}{1 Why quantum semantics?}
      \begin{center}
\begin{tikzpicture}[x=1cm,y=1cm,>=Latex]
  \tikzset{
    term/.style={align=center,text=navy,font=\fontsize{24}{28}\selectfont},
    edge/.style={line width=3pt,color=teal},
    oppose/.style={line width=3pt,color=gray!70,dashed},
    tag/.style={fill=white,inner sep=4pt,text=teal,font=\fontsize{24}{28}\selectfont\bfseries}
  }
  \node[term] (s1) at (0,4.8) {$S_1$\\\emph{female}};
  \node[term] (s2) at (8,4.8) {$S_2$\\\emph{male}};
  \node[term] (ns2) at (0,0) {$\neg S_2$\\\emph{non-male}};
  \node[term] (ns1) at (8,0) {$\neg S_1$\\\emph{non-female}};
  \draw[edge] (s1) -- node[tag,above] {$\sigma_z$} (s2);
  \draw[oppose] (ns2) -- (ns1);
  \draw[edge,->] (s1) -- node[tag,left] {$\sigma_y$} (ns2);
  \draw[edge,->] (s2) -- node[tag,right] {$\sigma_y$} (ns1);
  \draw[edge,<->] (s1) -- (ns1);
  \draw[edge,<->] (s2) -- (ns2);
  \node[tag] at (4,2.25) {$\sigma_x$};
  \node[font=\fontsize{20}{24}\selectfont,text=gray!75!black,fill=white,inner sep=3pt] at (4,-0.65) {subcontraries};
\end{tikzpicture}
\end{center}

      Pauli operators encode relations; context hides antonyms (Galofaro et al., 2018; Dubois \& Toffano, 2026). Composition: Coecke et al. (2010).
    \end{block}
  \end{column}
  \begin{column}{\sepwidth}\end{column}
  \begin{column}{\widecolwidth}
    \begin{block}{2 Three models, one rule}
      \begin{center}
\begin{tikzpicture}[x=1cm,y=1cm,>=Latex]
  \tikzset{
    labelbox/.style={draw=navy,fill=navy,text=white,rounded corners=5pt,minimum height=1.55cm,text width=5.1cm,align=center,font=\fontsize{24}{28}\selectfont\bfseries,inner sep=7pt},
    classical/.style={draw=burgundy!75,fill=pale,rounded corners=5pt,minimum height=1.55cm,text width=8.8cm,align=center,font=\fontsize{24}{28}\selectfont,inner sep=7pt},
    quantum/.style={draw=teal,fill=teal!12,rounded corners=5pt,minimum height=1.55cm,text width=13.6cm,align=center,font=\fontsize{24}{28}\selectfont,inner sep=7pt},
    readout/.style={draw=navy,fill=white,rounded corners=5pt,minimum height=1.55cm,text width=10.0cm,align=center,font=\fontsize{24}{28}\selectfont,inner sep=7pt},
    link/.style={draw=navy,line width=3pt,->}
  }
  \node[labelbox] (q0) at (0,4.6) {QCSE};
  \node[classical] (q1) at (9.2,4.6) {context};
  \node[quantum] (q2) at (22.7,4.6) {quantum ansatz};
  \node[readout] (q3) at (38.3,4.6) {bit word};
  \node[labelbox] (h0) at (0,2.3) {Hybrid};
  \node[classical] (h1) at (9.2,2.3) {embedding};
  \node[quantum] (h2) at (22.7,2.3) {8 qubits};
  \node[readout] (h3) at (38.3,2.3) {word head};
  \node[labelbox] (m0) at (0,0) {Hamiltonian};
  \node[classical] (m1) at (9.2,0) {count prior};
  \node[quantum] (m2) at (22.7,0) {$H$ evolution};
  \node[readout] (m3) at (38.3,0) {Born readout};
  \draw[link] (q0) -- (q1); \draw[link] (q1) -- (q2); \draw[link] (q2) -- (q3);
  \draw[link] (h0) -- (h1); \draw[link] (h1) -- (h2); \draw[link] (h2) -- (h3);
  \draw[link] (m0) -- (m1); \draw[link] (m1) -- (m2); \draw[link] (m2) -- (m3);
\end{tikzpicture}
\end{center}
\vspace{0.25cm}
\begin{center}
\fontsize{24}{29}\selectfont
\renewcommand{\arraystretch}{1.10}
\begin{tabularx}{0.98\linewidth}{@{}>{\bfseries}p{0.18\linewidth}p{0.25\linewidth}X@{}}
\toprule
Model & Size & Classical control\\
\midrule
QCSE & 10 audit; 14 full & constant prior\\
Hybrid & 8 qubits & identity, MLP, BERT\\
Hamiltonian & 10 qubits & dephased, MLP\\
\bottomrule
\end{tabularx}
\end{center}

      \vspace{0.25cm}
      \begin{center}
        \textbf{Same inputs, targets and budget; checkpoints selected on validation.}
      \end{center}
      \vspace{0.15cm}
      QCSE reimplementation: Varmantchaonala et al. (2026); MLP: Bengio et al. (2003).
    \end{block}
  \end{column}
\end{columns}

\vspace{1cm plus 1fill}

\begin{alertblock}{3 Results: fair controls change the story}
\vspace{0.4cm}
\begin{columns}[t,totalwidth=\textwidth]
  \begin{column}{\colwidth}
    \vspace{0.35cm}
    \textbf{A constant guess beats the published metric}\par
    \vspace{0.25cm}
    \begin{center}
\begin{tikzpicture}
\begin{axis}[
  ybar,
  width=0.88\linewidth,
  height=20cm,
  xmin=0.3,xmax=5.7,
  ymin=0,ymax=108,
  ylabel={Score (\%)},
  xtick={1.5,4.5},
  xticklabels={Half-bit,Exact word},
  ytick={0,50,100},
  bar width=1.2cm,
  enlarge x limits=false,
  legend style={at={(0.5,-0.19)},anchor=north,legend columns=2,column sep=0.8cm},
  every axis plot/.append style={line width=3pt},
  nodes near coords={\pgfmathprintnumber[fixed,precision=1]\pgfplotspointmeta},
  nodes near coords style={font=\fontsize{24}{28}\selectfont,anchor=south}
]
  \addplot[fill=teal,draw=teal] coordinates {(1,56.1) (4,0)};
  \addlegendentry{QCSE}
  \addplot[fill=gray!65,draw=gray!65] coordinates {(2,96.5) (5,22.8)};
  \addlegendentry{Constant prior}
  \draw[gray!75,densely dashed,line width=3pt] (axis cs:0.55,50) -- (axis cs:2.45,50);
  \node[font=\fontsize{24}{28}\selectfont,anchor=west,text=gray!70!black] at (axis cs:2.5,51) {50\% bit chance};
\end{axis}
\end{tikzpicture}
\end{center}
\par\smallskip
\begin{minipage}{\linewidth}
\RaggedRight\fontsize{24}{29}\selectfont
57 targets; 96.5\% is not word recovery.
\end{minipage}

  \end{column}
  \begin{column}{\sepwidth}\end{column}
  \begin{column}{\colwidth}
    \vspace{0.35cm}
    \textbf{Removing the circuit changes almost nothing}\par
    \vspace{0.25cm}
    \begin{center}
\begin{tikzpicture}
\begin{axis}[
  width=0.68\linewidth,
  height=18cm,
  xmin=0,xmax=60,
  ymin=0.55,ymax=2.45,
  xlabel={Top-1 accuracy (\%)},
  ytick={2,1},
  yticklabels={Left-only (L4),Bidirectional (L2R2)},
  xtick={0,30,60},
  legend style={at={(0.5,-0.28)},anchor=north,legend columns=3,column sep=0.55cm},
  xmajorgrids=true,
  ymajorgrids=false
]
  \addplot[only marks,color=teal,mark=*,mark size=7pt] coordinates {(33.32,2) (50.38,1)};
  \addlegendentry{Trained circuit}
  \addplot[only marks,color=gray!70!black,mark=square*,mark size=7pt] coordinates {(33.36,2) (50.03,1)};
  \addlegendentry{Identity}
  \addplot[only marks,color=burgundy,mark=diamond*,mark size=8pt] coordinates {(35.66,2) (54.89,1)};
  \addlegendentry{MLP}
\end{axis}
\end{tikzpicture}
\end{center}
\par\smallskip
\begin{minipage}{\linewidth}
\RaggedRight\fontsize{24}{29}\selectfont
L4 circuit/identity/MLP: 33.32/33.36/35.66\%. L2R2: 50.38/50.03/54.89\%. Across 15 runs: +0.23 point over identity at 2.7$\times$ compute; fold signs vary.
\end{minipage}

  \end{column}
  \begin{column}{\sepwidth}\end{column}
  \begin{column}{\colwidth}
    \vspace{0.35cm}
    \textbf{It learns, but the MLP stays ahead}\par
    \vspace{0.25cm}
    \begin{center}
\begin{tikzpicture}
\begin{axis}[
  width=0.84\linewidth,
  height=18cm,
  xmin=0.8,xmax=4.35,
  ymin=2.0,ymax=4.55,
  ylabel={Validation NLL (nats; lower is better)},
  xtick={1,2,3,4},
  xticklabels={First BPE,Product state,Count prior,$64\times$ pool},
  xticklabel style={font=\fontsize{24}{29}\selectfont,rotate=20,anchor=north east},
  ytick={2,2.5,3,3.5,4,4.5},
  legend style={at={(0.5,-0.24)},anchor=north,legend columns=2,column sep=0.5cm},
  every axis plot/.append style={line width=3pt}
]
  \addplot[color=teal,mark=*,line width=3pt] coordinates {(1,4.2773) (2,3.1975) (3,2.66776) (4,2.46546)};
  \addlegendentry{Hamiltonian}
  \addplot[color=burgundy,mark=square*,line width=3pt] coordinates {(1,3.4424) (2,3.2306) (3,2.41207) (4,2.22277)};
  \addlegendentry{MLP}
  \addplot[color=gray!75,densely dotted,line width=3pt,domain=1:4,samples=2] {3.0593};
  \addlegendentry{Bigram counts}
  \node[font=\fontsize{24}{28}\selectfont,anchor=west,text=navy] at (axis cs:3.55,2.48) {0.24 nats};
\end{axis}
\end{tikzpicture}
\end{center}
\par\smallskip
\begin{minipage}{\linewidth}
\RaggedRight\fontsize{24}{29}\selectfont
Validation milestones; recipes vary. Tatoeba: coherent $H$ 6.76\% vs dephased 5.71\%. Prompt “The history of”: “its first game ream, butvarBan7…”
\end{minipage}

  \end{column}
\end{columns}
\end{alertblock}

\vspace{1cm plus 1fill}

\begin{columns}[t,totalwidth=\textwidth]
  \begin{column}{\colwidth}
    \begin{block}{4 What it means}
      Gains came from embeddings, count priors and pooling.\\[0.35em]
      Limits: exact simulation, few seeds; final test unscored.\\[0.35em]
      Next: test relations directly.
    \end{block}
  \end{column}
  \begin{column}{\sepwidth}\end{column}
  \begin{column}{\widecolwidth}
    \begin{alertblock}{5 Take-home message}
      \fontsize{34}{41}\selectfont
      No advantage is demonstrated: matched classical models equal or beat every quantum or hybrid approach.\\[0.25em]
      A constant prior scored 96.5\% on the published criterion.\\[0.25em]
      Run identity controls to measure the circuit's contribution.
    \end{alertblock}
  \end{column}
\end{columns}

\vspace{1cm plus 1fill}

\begin{columns}[t,totalwidth=\textwidth]
  \begin{column}{\widecolwidth}
    \begin{block}{References}
      \fontsize{20}{24}\selectfont
      (Galofaro et al., 2018); (Dubois \& Toffano, 2026); (Coecke et al., 2010); (Varmantchaonala et al., 2026); (Bengio et al., 2003).
    \end{block}
  \end{column}
  \begin{column}{\sepwidth}\end{column}
  \begin{column}{\colwidth}
    \begin{block}{Credits}
      \fontsize{20}{24}\selectfont
      AI and OneHealth, 20 Nov 2026. LLMs aided literature, code and drafting; results checked.
    \end{block}
  \end{column}
\end{columns}
\end{frame}
\end{document}
